\section{Measurement Selection and Precision Planning}

\subsection{Selection of independent neuropsychological outcomes}
The transfer endpoint should measure an executive-control construct without reproducing the practiced mobile response mapping. Candidate domains include response inhibition and cognitive flexibility. Selection will follow an explicit evidence review of age coverage, test--retest reliability, normative applicability to the recruited population, susceptibility to practice effects, administration time, and availability of alternate forms. A commercial or named test is not assumed valid solely because it is commonly used. The primary score, direction of improvement, scoring algorithm, and assessor training requirements will be frozen before recruitment.

Construct overlap must be distinguished from item overlap. A task can target inhibition while using unfamiliar stimuli and response modalities, whereas a nominally different task may reproduce the same trained rule. The protocol will document this distinction and the expected near-transfer or farther-transfer distance. Parent reports of everyday functioning can be secondary external-validity measures but should not be combined with task performance into an unvalidated composite.

\subsection{Precision of the longitudinal transfer contrast}
For a two-arm randomized trial with a continuous post-training score, equal allocation, common residual standard deviation \(\sigma_e\), and an approximate normal confidence interval, a planning approximation for the confidence-interval half-width \(h\) is
\[
h \approx z_{1-\alpha/2}\,\sigma_e\sqrt{\frac{4}{N}},
\]
where \(N\) is the total number of participants with analyzable outcomes. Hence,
\[
N \approx \frac{4z_{1-\alpha/2}^{\,2}\sigma_e^2}{h^2}.
\]
These equations illustrate precision planning, not a justified sample size for the present study. Residual variability must be estimated from relevant pilot data or external evidence. If baseline adjustment is anticipated, the residual variance rather than the unadjusted endpoint variance should be used, with its assumptions stated. Unequal allocation, clustered recruitment, covariate imbalance, and non-Gaussian outcomes require an appropriate revised calculation.

Anticipated outcome attrition may be incorporated through a planning inflation \(N_{\mathrm{recruit}}\approx N/(1-r)\), where \(r\) is the expected missing-outcome proportion. This adjustment does not correct informative missingness and cannot replace sensitivity analysis. Precision planning should also evaluate the probability of excluding clinically meaningful harm, not merely the chance of obtaining a small \(p\)-value.

\subsection{Repeated-measure crossover information}
In the three-context crossover, effective information depends on the number of children, within-child correlation, between-child variability, order effects, and technically valid blocks. Increasing blocks for the same child does not substitute indefinitely for recruiting independent children. Simulation-based planning will reproduce the intended condition sequence, block-level error structure, incomplete AR tracking, and child-level dropout. The simulated decision criteria will include confidence-interval coverage for E--S and A--E and the probability of an incorrect non-inferiority conclusion under unfavorable scenarios.

\subsection{Pilot-derived assumptions and sensitivity grid}
The feasibility stage will estimate usable-trial fractions, baseline error prevalence, latency variability, device qualification rate, within-child correlation, and tolerability-related discontinuation. Planning scenarios will vary these quantities across credible ranges rather than rely on a single optimistic estimate. The final recruitment target, non-inferiority margin, transfer threshold, alpha allocation, and stopping rules must be justified and registered before confirmatory data collection.

\subsection{Interpretation safeguards}
A statistically precise gain in movement exposure is not proof of improved regulation. Failure to reject a null difference in cognitive accuracy is not evidence of non-inferiority. An independent transfer gain cannot be attributed specifically to AR unless the randomized comparison isolates that component; a varied-context versus fixed-context trial identifies the training-schedule package. Exploratory associations between movement, engagement, and transfer do not establish mediation without additional assumptions and a dedicated analysis.